\documentclass[10pt,a4paper]{amsart}
\usepackage[margin=19mm]{geometry}
\usepackage{amsmath,amssymb,mathtools}
\usepackage[scr=rsfs,cal=euler]{mathalfa}
\usepackage{xfrac}
\usepackage[unicode,hidelinks,bookmarks=false]{hyperref}
\newcommand{\ie}{i.e.\ }
\newcommand{\cf}{cf.\ }
\newcommand{\resp}{respectively\ }
\newcommand{\Subsection}{\subsection}
\providecommand{\nobreakdash}{\nobreak\hbox{-}}
\setlength{\emergencystretch}{2em}
\multlinegap0pt
\usepackage{amssymb}
\usepackage{xcolor}
\newtheorem{proposition}{Proposition}
\newtheorem{lemma}{Lemma}
\newtheorem{theorem}{Theorem}
\title{Singularity Models of Finite-Time K\"ahler-Ricci Flows}
\author{Frederick Tsz-Ho Fong, Hung Tran}
\date{Source: arXiv:2609.03332v1 (CC BY 4.0)}
\begin{document}
\maketitle

\noindent\emph{Source and licence.} Singularity Models of Finite-Time K\"ahler-Ricci Flows. Version arXiv:2609.03332v1 (CC BY 4.0), distributed by arXiv under the Creative Commons Attribution 4.0 International licence, http://creativecommons.org/licenses/by/4.0/, as recorded in the arXiv licence metadata for this exact version and shown on the abstract page https://arxiv.org/abs/2609.03332v1. Full licence notice: \texttt{licenses/arxiv-2609.03332-CC-BY-4.0.txt}.\par\smallskip

\noindent\textit{Editorial scope.} Counted unit: the article's theorem thm:h2pp-upper-bound (source0.tex lines 1546--1556). The article's complete original proof of it is delivered with it (source0.tex lines 1558--1623, one \texttt{proof} environment of 66 lines). No mathematics is written, repaired or re-derived: the statement and the proof are the article's own lines, in the article's own notation, and no retained line is substituted or re-flowed.  Retained uncounted as the source-local context the proof needs: lines 917--919; lines 1159--1161; lines 1477--1483; lines 1519--1523.  Nothing was omitted from any retained span, so the package carries no omission record.  No retained line contains a citation, so the wrapper carries no bibliography and no bibliography entry.  No retained line contains a figure environment, an image-inclusion command or drawing code, so no image is delivered and the package declares no asset dependency.  Editorial formatting only, in addition: an AMS \texttt{amsart} preamble that restates the article's own declarations and macros used by the retained lines, namely the environments \texttt{proposition}, \texttt{lemma}, \texttt{theorem} on separate counters, together with the standard packages those lines need.  The retained environments print in this document as Proposition 1, Lemma 1, Theorem 1 in order of appearance, whereas the article prints the same items with the numbering of its own full document.  The complete native source of the article as distributed in the arXiv e-print for this exact version is bundled unchanged as \texttt{sources/209288/source0.tex}.
		\begin{equation}\label{eq:li-yau}
		\sup_{\widehat{M}} \frac{| \nabla u |^2}{u} (\cdot, t) \;\leq\; C_0 \;:=\; \max \left\{ \sup_{\widehat{M}} \frac{| \nabla u |^2}{u} (\cdot, 0) \,, \; 4\, p_i \right\} \qquad \text{for all } t \in [0, T) \,.
	\end{equation}

	\begin{equation}\label{eq:linear-lower-bound}
		f_k^2 (\cdot, t) \;\geq\; c\, \big( T - t \big) \qquad \text{on } \widehat{M} \times [0, T) \,, \quad \text{for every } k \geq 2 \,,
	\end{equation}

\begin{proposition}[the flow equation for $(h^2)''$]\label{prop:h2pp-evolution}
	On $\{ 0 \leq f_1^2 \leq L (t) ,\; 0 \leq t < T \}$,
	\begin{align}\label{eq:h2pp-evolution}
		\partial_t (h^2)'' &=\Delta (h^2)'' - ((h^2)'')^2 + R,\\
		R :&= -m \partial^2_{f_1^2} [ h^4/f_1^4 ] - \sum_{k \ge 2} n_k q_k^2 \partial^2_{f_1^2} [ h^4/f_k^4 ].\nonumber
	\end{align}
\end{proposition}

\begin{lemma} \label{lem:residual-source-bound}
	We have, where $(h^2)'' \ge 0$,
	\[R_{res} := -\sum_{k \ge 2} n_k q_k^2 \partial^2_{f_1^2} [ h^4/f_k^4 ] \le \frac{\beta}{(T-t)^2}, \quad \text{where} \quad \beta = \frac{2C_0^2 n'}{c^2}.\]
	Here $C_0$ is the Li--Yau constant of \eqref{eq:li-yau}, $c$ that of \eqref{eq:linear-lower-bound}, and $n' = \sum_{k \geq 2} n_k$.
\end{lemma}

\begin{theorem}[the negative part of the fiber curvature is Type I]\label{thm:h2pp-upper-bound}
	Let $C_0$ be the Li--Yau constant of \eqref{eq:li-yau}, $c$ that of \eqref{eq:linear-lower-bound}, $n' = \sum_{k \geq 2} n_k$, and set
	\[
	\gamma \;:=\; \max \Big\{ \tfrac{1}{2} \Big( 1 + \sqrt{ 1 + 8\, C_0^2\, n' / c^2 } \Big) \,,\;\; T \cdot \max_{\widehat{M}} ( h^2 )'' ( \cdot, 0 ) \Big\} \,.
	\]
	Then
	\begin{equation}\label{eq:h2pp-upper-bound}
		( h^2 )'' \;\leq\; \frac{\gamma}{T - t} \qquad \text{on } \widehat{M} \times [ 0, T ) \,.
	\end{equation}
	
\end{theorem}

\begin{proof}
	Our goal is to show that, at a spacial maximum of $(h^2)''$, for $R_1 := -m \partial^2_{f_1^2}[ h^4/f_1^4 ]$
	\[\Delta(h^2)''+R_1\leq 0. \]
	By the maximum principle, immediately, at such a point, $\Delta(h^2)''\leq 0$. Now we investigate the term $R_1$. By direct calculation,
	\[
	R_1 \;=\; \frac{2 m}{f_1^2} \cdot \frac{h^2}{f_1^2} \left( 2 \left( \frac{h^2}{f_1^2} \right)' - ( h^2 )'' \right) \;-\; 2 m \left( \left( \frac{h^2}{f_1^2} \right)' \right)^2 \,.
	\]
 
	
	Near the nut stratum $Q_0$ where $f_1 \to 0$, this term contains a potential $1/f_1^8$ coordinate singularity. To bypass this, we work directly with the ratio $h^2/f_1^2$ and the integral identity:
	\begin{equation}
		\left( \frac{h^2}{f_1^2} \right)' = \frac{1}{f_1^4} \int_0^{f_1^2} \zeta (h^2)''(\zeta) \, d\zeta,
	\end{equation}
%At any interior spatial maximum of $(h^2)''$, the integral identity forces the slope of the gradient to satisfy $2(h^2/f_1^2)' \le (h^2)''$, which algebraically guarantees that $R_1 \le 0$.
 If the spatial maximum of $(h^2)''$ happens at an interior point  $\zeta_0$, the integrand is bounded above by the maximum non-negative value $(h^2)''(\zeta_0)$: 
 
 $$\left( \frac{h^2}{f_1^2} \right)'(\zeta_0) \le \frac{1}{\zeta_0^4} (h^2)''(\zeta_0) \int_0^{\zeta_0} \zeta \, d\zeta = \frac{1}{\zeta_0^4} (h^2)''(\zeta_0) \frac{\zeta_0^2}{2} = \frac{1}{2\zeta_0^2} (h^2)''(\zeta_0)$$
 Since $\zeta_0^2 = f_1^4$, we multiply both sides by $2 f_1^2$ to get: $$2 \left( \frac{h^2}{f_1^2} \right)'(\zeta_0) \le (h^2)''(\zeta_0)$$
 Thus, $R_1\leq 0$. \\
 
If the spatial maximum of $(h^2)''$ is on the bolt $Q_\ell$, we have $h^2/f_1^2 = 0$ there. So $R_1 = - 2 m \big( ( h^2 / f_1^2 )' \big)^2 \leq 0$.\\

	At the nut stratum $Q_0$ (where $f_1^2 = 0$), the interior integral identity does not apply directly. Instead, we expand the Laplacian and the collapsing source term together. %more care is needed, because the interior argument genuinely fails there: in the formula for $R_1$ the prefactor blows up, $\tfrac{2 m}{f_1^2} \cdot \tfrac{h^2}{f_1^2} \to \infty$ as $f_1^2 \to 0$ (since $h^2 / f_1^2 \to 2$ at $Q_0$), while the bracket $2 ( h^2 / f_1^2 )' - ( h^2 )''$ tends to $0$ --- an $\infty \cdot 0$ competition that must be resolved by expanding one order deeper; and the integral inequality above used $\zeta_0 > 0$. 
	As seen in the proof of Prop. \ref{prop:h2pp-evolution}, near $Q_0$ the function $( h^2 )''$ is a smooth function of $f_1^2$, so it admits the one-sided expansion:	
	\[
	( h^2 )'' \;=\; y (0) \;+\; \big( ( h^2 )'' \big)' ( 0^+ )\, f_1^2 \;+\; O \big( ( f_1^2 )^2 \big) \,, \qquad y (0) \;:=\; ( h^2 )'' \big|_{Q_0} \,.
	\]
	The value $y (0)$ is the spatial maximum $y (t)$ in the present case. Integrating twice from the boundary data $h^2 = 0$, $( h^2 )' = 2$ at $f_1^2 = 0$,
	\[
	h^2 \;=\; 2 f_1^2 \;+\; \frac{y (0)}{2}\, ( f_1^2 )^2 \;+\; \frac{1}{6}\, \big( ( h^2 )'' \big)' ( 0^+ )\, ( f_1^2 )^3 \;+\; O \big( ( f_1^2 )^4 \big) \,.
	\]
	Therefore, we have	
	\[
	\left( \frac{h^2}{f_1^2} \right)' \;=\; \frac{y (0)}{2} \;+\; \frac{1}{3}\, \big( ( h^2 )'' \big)' ( 0^+ )\, f_1^2 \;+\; O \big( ( f_1^2 )^2 \big) \,, \qquad \text{and} \quad \left( \frac{h^2}{f_1^2} \right)' (0) \;=\; \frac{y (0)}{2} \,.
	\]	
	Consequently,	
	\[
	2 \left( \frac{h^2}{f_1^2} \right)' - ( h^2 )'' \;=\; - \frac{1}{3}\, \big( ( h^2 )'' \big)' ( 0^+ )\, f_1^2 \;+\; O \big( ( f_1^2 )^2 \big) \,.
	\]
	That is, the bracket in the formula of $R_1$ thus vanishes to {first order} in $f_1^2$: exactly enough to absorb the singular prefactor $\tfrac{2 m}{f_1^2} \cdot \tfrac{h^2}{f_1^2} = \tfrac{4 m}{f_1^2} \big( 1 + O ( f_1^2 ) \big)$. Hence, $R_1$ extends continuously to $Q_0$ with value
	
	\[
	R_1 \big|_{Q_0} \;=\; - \frac{4 m}{3}\, \big( ( h^2 )'' \big)' ( 0^+ ) \;-\; \frac{m}{2}\, y (0)^2 \,.
	\]
	
	%Unlike in the interior and at $Q_\ell$, this need \textbf{not} be nonpositive on its own: maximality of $( h^2 )''$ at $f_1^2 = 0$ forces $\big( ( h^2 )'' \big)' ( 0^+ ) \leq 0$, so the first term is $\geq 0$. 
	
	The rescue is the Laplacian, which at $Q_0$ is not merely non-positive but carries the quantitative stratum drift of Prop \ref{prop:h2pp-evolution}. That is, $\Delta \big( ( h^2 )'' \big) = h^2\, ( h^2 )'''' + ( h^2 )'''\, \Delta ( f_1^2 )$; the first term vanishes at $Q_0$ (where $h^2 = 0$) and $\Delta ( f_1^2 ) |_{Q_0} = 2 ( m + 1 )$, so
	
	\[
	\Delta \big( ( h^2 )'' \big) \Big|_{Q_0} \;=\; 2 ( m + 1 )\, \big( ( h^2 )'' \big)' ( 0^+ ) \,.
	\]
	Thus, at a spacial maximum of $(h^2)''$, for $R_1 := -m \partial^2_{f_1^2}[ h^4/f_1^4 ]$
	\[\Delta(h^2)''+R_1\leq 0. \]
	Let $y(t) := \max_{\widehat{M}} (h^2)''(\cdot, t) \ge 0$ then  $y (t)$ is locally Lipschitz. Applying the estimate above in conjunction with Lemma \ref{lem:residual-source-bound}
	%the maximum principle to our regularized, self-damping PDE, the Laplacian and the collapsing source terms drop out, leaving 
	yields the Riccati differential inequality, in the sense of Dini derivatives,
	\begin{equation*}
		\frac{dy}{dt} \le -y^2 + \frac{\beta}{(T-t)^2}.
	\end{equation*}
	By standard ODE comparison, since the initial metric satisfies the boundary condition $y(0) \le \gamma / T$, we conclude:
	\begin{equation}
		y(t) \le \frac{\gamma}{T - t}
	\end{equation}
	for all $t \in [0, T)$. This completes the proof.
\end{proof}

\end{document}
